% SPDX-License-Identifier: GPL-3.0-or-later
% Copyright (C) 2026 Alessandro Greco
%% Testo alternativo per l'accessibilita' (usato solo nell'EPUB).
\altfig{Cinque compiti crittografici in Tor: riservatezza con AES-CTR nelle celle e AES-GCM o ChaCha20-Poly1305 nei canali TLS (capitoli 10 e 7); integrità con SHA-1 e SHA3-256 per i digest e POLYVAL per i tag CGO (capitolo 10); accordo sulle chiavi tramite X25519 in ntor e ntor-v3 (capitolo 9); autenticazione con Ed25519 e RSA legacy per certificati e consenso (capitoli 7 e 12); derivazione delle chiavi con HKDF-SHA256 e SHAKE-256 (capitolo 9).}
% I cinque compiti della crittografia in Tor e le primitive usate
\begin{tikzpicture}[font=\sffamily\scriptsize,
  compito/.style={draw=torviola,fill=torviola!8,rounded corners=2pt,minimum width=2.6cm,minimum height=0.7cm},
  prim/.style={draw=torverde!60!black,fill=torverde!12,rounded corners=2pt,text width=3.9cm,align=center,minimum height=0.7cm},
  dove/.style={font=\sffamily\tiny,text=torgrigio,anchor=west}]
  \foreach \i/\c/\p/\d in {%
    0/Riservatezza/{AES-CTR; TLS con AES-GCM o ChaCha20-Poly1305}/{celle (cap.~10), canali (cap.~7)},
    1/Integrità/{SHA-1 e SHA3-256, POLYVAL}/{digest delle celle RELAY, tag di autenticazione CGO (cap.~10)},
    2/{Accordo sulle chiavi}/{X25519 (Diffie-Hellman)}/{ntor, ntor-v3 (cap.~9)},
    3/Autenticazione/{Ed25519, RSA (in uscita)}/{certificati, consenso (cap.~7, 12)},
    4/Derivazione/{HKDF-SHA256, SHAKE-256}/{chiavi di circuito (cap.~9)}} {
    \node[compito] (c\i) at (0,-\i*1.0) {\c};
    \node[prim] (p\i) at (4.4,-\i*1.0) {\p};
    \draw[->,torgrigio] (c\i) -- (p\i);
    \node[dove] at (6.8,-\i*1.0) {\d};
  }
\end{tikzpicture}
